\section{Knowledge Distillation for Language Models}

\subsection{Classical Off-Policy Logit Distillation}

Knowledge distillation (KD) has established itself as a foundational technique for model compression and efficiency in natural language processing \parencite{hinton2015distillingknowledgeneuralnetwork}. The classical framework, as formalized by Hinton et al., seeks to transfer predictive capabilities from a large teacher model to a smaller student model by minimizing the divergence between their output probability distributions. Specifically, the student is trained to match the soft targets (logits) produced by the teacher, which are scaled by a temperature parameter $T > 1$ to smooth the distribution:
\begin{equation}
    \mathcal{L}_{KD} = T^2 \cdot D_{KL}\left(\sigma\left(\frac{\mathbf{z}_T}{T}\right) \parallel \sigma\left(\frac{\mathbf{z}_S}{T}\right)\right),
\end{equation}
where $\mathbf{z}_T$ and $\mathbf{z}_S$ represent the logits of the teacher and student models, respectively, and $\sigma$ denotes the softmax function. These soft targets carry "dark knowledge" --- subtle statistical correlations between classes or tokens that hard labels fail to capture. 

In the domain of language modeling, this classical formulation has been widely applied to compress representation models, as exemplified by DistilBERT \parencite{sanh2020distilbertdistilledversionbert}. However, classical logit distillation operates under an *off-policy* paradigm. During training, the teacher's logits are computed using reference text (gold-standard tokens or teacher-generated prefixes). Consequently, the student is only supervised on trajectories generated by an external policy, creating a critical misalignment during auto-regressive inference.

\subsection{Exposure Bias in Sequential Reasoning}

The primary bottleneck of off-policy logit distillation in generative tasks is the phenomenon of exposure bias. Exposure bias arises from the discrepancy in the prefix distribution between the training phase and the inference phase. During off-policy training (often implemented via teacher-forcing), the student model is exposed exclusively to gold-standard historical tokens as prefixes. During inference, however, the student must generate text auto-regressively, conditioning its predictions on its own previously generated tokens, which may contain errors.

As a consequence of this training-inference discrepancy, any small error made by the student at the beginning of a sequence alters its input distribution away from the training distribution. Since the student has never been trained to recover from its own mistakes, these errors compound and propagate exponentially along the generation trajectory. In the context of legal reasoning (e.g., Chain-of-Thought reasoning or syllogistic deduction), logical consistency is paramount. A single logical misstep or vocabulary error in the initial reasoning steps will inevitably derail the entire deduction process, leading to flawed legal conclusions. Off-policy distillation is therefore inherently limited in its ability to produce small language models capable of robust sequential reasoning.

\subsection{On-Policy Distillation}

To mitigate the effects of exposure bias, recent research has shifted toward *on-policy distillation* frameworks \parencite{agarwal2024onpolicydistillationlanguagemodels}. Under this paradigm, the student model generates its own sequence rollouts (trajectories) during the training phase. The teacher model then acts as an evaluator, providing token-level supervision (logits or feedback) on the student's self-generated prefixes:
\begin{equation}
    \mathcal{L}_{\text{on-policy}} = \mathbb{E}_{y \sim \pi_S(\cdot|x)} \left[ D_{KL}\left(\pi_T(\cdot|[x, y]) \parallel \pi_S(\cdot|[x, y])\right) \right],
\end{equation}
where $y$ is a sequence sampled from the student policy $\pi_S$ given the input prompt $x$, and $\pi_T$ is the teacher policy.

By utilizing student-generated rollouts, on-policy distillation exposes the student to its own generation distribution during training. The teacher model's supervision on these paths acts as a corrective feedback loop, teaching the student how to navigate and recover from its own mistakes. This iterative correction makes on-policy distillation highly effective for complex reasoning domains, such as mathematics \parencite{adarsh-etal-2025-siked} and structured legal reasoning, where the precision of sequential trajectories dictates the validity of the final outcome \parencite{song2026surveyonpolicydistillationlarge, zhao2026selfdistilledreasoneronpolicyselfdistillation}.

\subsection{Top-K Sparse Logits vs. Full-Vocabulary Distillation}

Implementing on-policy distillation introduces significant computational challenges, primarily centered on logit transmission and memory utilization. In *full-vocabulary distillation*, the Kullback-Leibler divergence is computed across the entire vocabulary space of the tokenizer (which typically ranges from 32,000 to over 150,000 tokens in modern models). While full-vocabulary distillation provides a rich and mathematically stable gradient signal, it requires the teacher model to execute a full forward pass and output the complete logit matrix for every student-generated token. This leads to substantial GPU memory overhead and slow training throughput.

Conversely, *top-K sparse logits* (or sampled-token) distillation limits the supervision signal to only the top $K$ tokens with the highest probabilities under the teacher's distribution (e.g., $K=50$). Although this sparse representation dramatically reduces memory usage and transmission bandwidth between the teacher and student processes, it introduces two distinct drawbacks. First, truncating the vocabulary distribution discards the long tail of the distribution, which contains valuable dark knowledge regarding negative constraints. Second, sparse logit matching can lead to high gradient variance during optimization, making the training process less stable compared to full-vocabulary matching. The trade-off between computational efficiency and distributional completeness remains a core design decision in the deployment of on-policy distillation pipelines.
